\chapter{Artificial Intelligence for Component Identification}
\label{ch:pfe-ai}

\section{Introduction}

Artificial intelligence can assist software-composition analysis when the
available evidence is incomplete, noisy, or too large to inspect manually. In
Android APK analysis, a model may help classify evidence, rank suspicious
cases, identify patterns linked to obfuscation, or support report triage. It
does not replace the need for explicit static evidence, reference databases,
and conservative security interpretation.

This chapter introduces the AI concepts used to understand such assistance.
It presents machine learning as a supervised classification problem, explains
feature representations for APK and DEX evidence, defines common evaluation
metrics, and describes deep-learning and Transformer models at a high level.
The chapter also clarifies the authority boundary of AI in an SBOM workflow:
models can help analysts interpret risk, but they should not silently change
component coordinates, versions, CycloneDX rows, or vulnerability labels.

\section{Machine Learning Concepts}

\subsection{Learning from Labeled Examples}

Machine learning is a family of methods that infer a decision rule from data
instead of relying only on manually written rules. In supervised learning, the
training set contains examples \(x_i\) and labels \(y_i\). The objective is to
learn a function \(f\) that predicts a label for a new example:

\begin{equation}
\label{eq:pfe-ai-supervised-learning}
\hat{y}=f(x).
\end{equation}

For component-identification support, an example may represent a candidate
match, a feature vector extracted from DEX evidence, or a row in an analysis
report. The label may indicate whether the candidate is strong, ambiguous,
unsupported, or suitable only for report-only investigation. The prediction is
useful only when the input features and labels are well defined.

\subsection{Classification}

Classification assigns each example to one class among a finite set of
possible labels. Binary classification uses two labels, such as positive and
negative. Multiclass classification uses more than two labels, such as
\texttt{accepted}, \texttt{ambiguous}, \texttt{risk-only}, and
\texttt{unsupported}. Figure~\ref{fig:pfe-ai-classification-flow} illustrates
the general supervised classification loop.

\begin{figure}[H]
\centering
\resizebox{0.92\textwidth}{!}{%
\begin{tikzpicture}[
  node distance=8mm and 9mm,
  box/.style={draw, rounded corners, align=center, text width=0.21\textwidth,
    minimum height=11mm, fill=blue!5},
  model/.style={draw, rounded corners, align=center, text width=0.20\textwidth,
    minimum height=11mm, fill=warningorange!14},
  outputnode/.style={draw, rounded corners, align=center, text width=0.21\textwidth,
    minimum height=11mm, fill=evidencegreen!10},
  arrow/.style={-{Latex[length=2.5mm]}, thick}
]
\node[box] (data) {Labeled examples};
\node[box, right=of data] (features) {Feature representation};
\node[model, right=of features] (train) {Learning algorithm};
\node[outputnode, right=of train] (model) {Trained classifier};
\node[box, below=of features] (new) {New APK/report evidence};
\node[outputnode, below=of model] (pred) {Predicted class};

\draw[arrow] (data) -- (features);
\draw[arrow] (features) -- (train);
\draw[arrow] (train) -- (model);
\draw[arrow] (new) -- (pred);
\draw[arrow] (model) -- (pred);
\end{tikzpicture}
}
\caption{General supervised classification workflow for APK evidence analysis.}
\label{fig:pfe-ai-classification-flow}
\end{figure}

The prediction is not automatically a security fact. It is a model output that
must be interpreted together with evidence provenance, validation data, error
rates, and the cost of incorrect decisions.

\subsection{Training, Validation, and Test Sets}

A dataset is usually separated into training, validation, and test sets. The
training set fits the model parameters. The validation set helps choose
hyperparameters and compare candidate configurations. The test set is kept
aside until the final evaluation so that the reported score better estimates
generalization to unseen data.

For APK-related evidence, splitting must avoid leakage. Examples derived from
the same application, component family, or near-duplicate report can be very
similar. If related rows appear in both training and test sets, the model may
memorize local patterns instead of learning a useful rule. Grouped splits are
therefore preferable when several rows come from the same application or
experiment.

\subsection{Class Imbalance}

Software-security datasets are often imbalanced. A common class may contain
many examples, while rare but important classes contain few. The class
distribution is imbalanced when the largest class dominates the smallest
class. For example, a dataset may contain 7,811 rows in its largest class and
594 rows in its smallest class. This imbalance motivates balanced class
weights during training and macro-averaged evaluation metrics during
assessment.

If imbalance is ignored, a classifier can obtain a high global accuracy by
favoring the majority class. Security analysis usually cares about minority
classes too, because rare unsupported, ambiguous, or risky cases may be the
ones that require analyst attention.

\section{Feature Representations}

\subsection{Structured Features}

Structured features convert an analysis object into numerical or categorical
values. In APK evidence analysis, such features may include match scores,
score margins, package overlap, class overlap, method overlap, literal
matches, API-reference evidence, version distance, source-validation status,
obfuscation indicators, or report labels. Table~\ref{tab:pfe-ai-feature-types}
summarizes common feature families.

\begin{table}[H]
\centering
\caption{Feature families used in AI-assisted APK evidence analysis.}
\label{tab:pfe-ai-feature-types}
\begin{tabularx}{\textwidth}{L{0.24\textwidth}X}
\toprule
\textbf{Feature family} & \textbf{Interpretation}\\
\midrule
Score features & Similarity values, confidence values, and margins between
top candidates.\\
Structural features & Class, method, dependency, call, and inheritance
patterns extracted from DEX evidence.\\
Lexical features & Package names, descriptors, literals, resource names, and
tokens preserved in the APK.\\
Context features & Application family, reference coverage, source-validation
status, transformation indicators, and experiment lane.\\
Report features & Existing report labels, uncertainty notes, risk flags, and
human-readable evidence summaries.\\
\bottomrule
\end{tabularx}
\end{table}

Structured features are interpretable and compatible with classical machine
learning models. Their limitation is that important information may be lost
when complex code evidence is reduced to a fixed set of columns.

\subsection{Token and Sequence Features}

Some evidence is naturally sequential. Smali tokens, opcode streams, method
descriptors, call chains, and textual report rows contain ordered symbols.
Tokenization converts these symbols into a vocabulary, and the model learns
from their order and co-occurrence. Figure~\ref{fig:pfe-ai-feature-pipeline}
shows the difference between structured and sequence representations.

\begin{figure}[H]
\centering
\resizebox{0.94\textwidth}{!}{%
\begin{tikzpicture}[
  node distance=8mm and 8mm,
  box/.style={draw, rounded corners, align=center, text width=0.22\textwidth,
    minimum height=11mm, fill=blue!5},
  model/.style={draw, rounded corners, align=center, text width=0.22\textwidth,
    minimum height=11mm, fill=warningorange!13},
  arrow/.style={-{Latex[length=2.5mm]}, thick}
]
\node[box] (apk) {APK or report evidence};
\node[box, right=of apk] (structured) {Structured vector};
\node[model, right=of structured] (classical) {Classical ML};
\node[box, below=of structured] (tokens) {Token sequence};
\node[model, right=of tokens] (deep) {Deep model};
\node[box, right=of classical] (decision) {Report-only prediction};

\draw[arrow] (apk) -- (structured);
\draw[arrow] (structured) -- (classical);
\draw[arrow] (apk) |- (tokens);
\draw[arrow] (tokens) -- (deep);
\draw[arrow] (classical) -- (decision);
\draw[arrow] (deep) -- (decision);
\end{tikzpicture}
}
\caption{Structured and sequence feature paths for AI-assisted evidence analysis.}
\label{fig:pfe-ai-feature-pipeline}
\end{figure}

Sequence models can capture patterns that are difficult to express manually,
but they require more data and are often harder to explain. In security
reporting, explainability matters because analysts need to know why a warning
or recommendation was produced.

\section{Classical Machine Learning}

\subsection{Decision Trees and Random Forests}

A decision tree predicts a label by applying a sequence of feature tests. It
is easy to inspect but can overfit training data. A Random Forest reduces this
risk by training many decision trees on different samples and feature subsets,
then combining their votes \cite{breiman2001randomforests}. If \(T_j(x)\) is
the prediction of tree \(j\), a forest with \(m\) trees predicts the majority
vote:

\begin{equation}
\label{eq:pfe-ai-random-forest-vote}
\hat{y}=\operatorname*{mode}_{j=1}^{m} T_j(x).
\end{equation}

\begin{figure}[H]
\centering
\resizebox{0.95\textwidth}{!}{%
\begin{tikzpicture}[
  node distance=7mm and 13mm,
  box/.style={draw, rounded corners, align=center, text width=0.22\textwidth,
    minimum height=10mm, fill=blue!5},
  tree/.style={draw, rounded corners, align=center, text width=0.18\textwidth,
    minimum height=10mm, fill=warningorange!14},
  vote/.style={draw, rounded corners, align=center, text width=0.22\textwidth,
    minimum height=10mm, fill=evidencegreen!10},
  arrow/.style={-{Latex[length=2.5mm]}, thick}
]
\node[box] (input) {Structured evidence\\feature vector};
\node[tree, right=18mm of input] (t2) {Tree 2\\feature subset};
\node[tree, above=of t2] (t1) {Tree 1\\bootstrap sample};
\node[tree, below=of t2] (t3) {Tree 3\\feature subset};
\node[vote, right=18mm of t2] (vote) {Majority vote\\or averaged probability};

\draw[arrow] (input.east) -- (t1.west);
\draw[arrow] (input.east) -- (t2.west);
\draw[arrow] (input.east) -- (t3.west);
\draw[arrow] (t1.east) -- (vote.west);
\draw[arrow] (t2.east) -- (vote.west);
\draw[arrow] (t3.east) -- (vote.west);
\end{tikzpicture}
}
\caption{Random Forest training and voting path for structured evidence.}
\label{fig:pfe-ai-random-forest}
\end{figure}

Random Forests are useful for tabular evidence because they handle mixed
feature types, nonlinear interactions, and feature importance analysis. Their
main limitation is that a high feature importance does not prove causality:
the model may rely on shortcuts that correlate with a label in one dataset but
do not generalize.

\subsection{Baselines}

A baseline is a simple method used as a comparison point. Common baselines
include majority-class prediction, simple threshold rules, or deterministic
heuristics. A complex model should be preferred only when it improves over
these baselines on held-out data and when the improvement is meaningful for
the intended decision. Baselines do two jobs here: they show whether the
learned model is genuinely extracting signal, and they expose when a more
complicated model is only learning a convenient shortcut from the dataset or
report format. In a security workflow, a weak baseline that already performs
well can be preferable to a larger model whose gain disappears outside the
training corpus.

\section{Deep Learning and Transformers}

\subsection{Neural Representations}

Deep learning models transform inputs through several layers of learned
representations. Instead of using only manually selected columns, they can
learn patterns from tokens, sequences, or embeddings. For APK analysis, this
can be useful when relevant evidence is distributed across many methods,
opcodes, or textual report fragments.

However, deep models usually need larger datasets than classical models. They
can also learn shortcuts, such as project-specific vocabulary or report
formatting artifacts \cite{geirhos2020shortcut}. Validation must therefore
check whether the model learns the intended evidence relation rather than an
accidental correlation. That is why the chapter treats deep models as
assistance layers: they can summarize or prioritize evidence, but they do not
replace the deterministic SBOM path.

\subsection{Transformer Encoders}

Transformers are neural sequence models based on attention
\cite{vaswani2017attention}. Attention allows the model to compare tokens
with other tokens in the same sequence and learn which relations are useful
for prediction. In simplified form, attention computes a weighted
combination of value vectors:

\begin{equation}
\label{eq:pfe-ai-attention}
\operatorname{Attention}(Q,K,V)=
\operatorname{softmax}\left(\frac{QK^{T}}{\sqrt{d_k}}\right)V.
\end{equation}

Here, \(Q\), \(K\), and \(V\) are query, key, and value matrices, and \(d_k\)
is the key dimension. In software-analysis tasks, tokens may represent
opcodes, descriptors, API calls, or structured report fields.

Transformers are powerful because the attention mechanism can compare
long-range tokens directly, which helps when a relevant signal is split across
multiple classes or report rows. They also demand caution: the same mechanism
can learn dataset shortcuts if the token stream contains coordinate names,
version labels, or report templates that correlate too strongly with the
target class. Figure~\ref{fig:pfe-ai-transformer-boundary} summarizes a
conservative Transformer use case. The model can recommend report-only
actions, but the authoritative SBOM path remains controlled by explicit
evidence and deterministic validation rules.

\begin{figure}[H]
\centering
\resizebox{0.82\textwidth}{!}{%
\begin{tikzpicture}[
  node distance=7mm,
  box/.style={draw, rounded corners, align=center, text width=0.64\textwidth,
    minimum height=9mm, fill=blue!5},
  ai/.style={draw, rounded corners, align=center, text width=0.64\textwidth,
    minimum height=9mm, fill=warningorange!14},
  safe/.style={draw, rounded corners, align=center, text width=0.64\textwidth,
    minimum height=9mm, fill=evidencegreen!10},
  arrow/.style={-{Latex[length=2.5mm]}, thick}
]
\node[box] (evidence) {Observed APK evidence and generated report rows};
\node[ai, below=of evidence] (model) {AI model suggests report-only triage};
\node[safe, below=of model] (guard) {Guardrail: no coordinate, version, SBOM, or CVE mutation};
\node[box, below=of guard] (analyst) {Analyst reviews recommendation with evidence};

\draw[arrow] (evidence) -- (model);
\draw[arrow] (model) -- (guard);
\draw[arrow] (guard) -- (analyst);
\end{tikzpicture}
}
\caption{Conservative authority boundary for Transformer-assisted report triage.}
\label{fig:pfe-ai-transformer-boundary}
\end{figure}

\section{Evaluation Metrics}

\subsection{Confusion Terms}

Classification performance is commonly described with true positives,
false positives, true negatives, and false negatives. For one class, precision
measures how many predicted positives are correct, while recall measures how
many true positives were recovered:

\begin{equation}
\label{eq:pfe-ai-precision}
\operatorname{Precision}=\frac{TP}{TP+FP}.
\end{equation}

\begin{equation}
\label{eq:pfe-ai-recall}
\operatorname{Recall}=\frac{TP}{TP+FN}.
\end{equation}

The \(F_1\) score combines precision and recall:

\begin{equation}
\label{eq:pfe-ai-f1}
F_1 = 2 \cdot \frac{\operatorname{Precision}\cdot\operatorname{Recall}}
{\operatorname{Precision}+\operatorname{Recall}}.
\end{equation}

Accuracy measures the proportion of all correct predictions:

\begin{equation}
\label{eq:pfe-ai-accuracy}
\operatorname{Accuracy}=\frac{TP+TN}{TP+TN+FP+FN}.
\end{equation}

\subsection{Macro and Weighted Averages}

In multiclass settings, metrics can be averaged in several ways. Macro
averaging computes the metric independently for each class and then averages
the class scores. It gives rare classes the same importance as common
classes. Weighted averaging also computes per-class scores but weights them
by class frequency. Accuracy can be misleading when the dataset is imbalanced,
so macro \(F_1\) is often more informative for security triage labels.

\section{Allowed and Forbidden AI Actions}

AI assistance in an SBOM workflow must have a restricted authority boundary.
The trained assistant may support three report actions:

\begin{itemize}
  \item suggest a report-only investigation path;
  \item suppress or weaken unsupported wording; and
  \item mark a row as risk-only for analyst attention.
\end{itemize}

It may not:

\begin{itemize}
  \item change the selected coordinate or version;
  \item add, remove, or modify CycloneDX components;
  \item change vulnerability identifiers or affected-version decisions; or
  \item promote report-only evidence into authoritative evaluation metrics.
\end{itemize}

This restriction is essential because model predictions are statistical
recommendations, not direct observations from the APK. A safe workflow keeps
AI outputs outside the authoritative component, version, SBOM, and CVE path
unless a separate validated process explicitly promotes them.

\section{Conclusion}

Artificial intelligence can support Android APK-to-SBOM analysis by learning
patterns in structured evidence, token sequences, and report rows. Classical
models such as Random Forests are suitable for interpretable tabular features,
while Transformer encoders can model ordered tokens and longer context.
However, both families require careful splitting, imbalance-aware metrics,
and validation against shortcut learning.

The most important principle is authority separation. AI can help prioritize
investigation, weaken unsupported language, or mark risk-only evidence for
review, but it should not silently alter coordinates, versions, SBOM
components, or vulnerability labels. The following chapter uses the state of
the art to position Android library-detection methods against these evidence
and authority requirements.
